\chapter{Флуктуационно-диссипативная мобильность общего следа}
\label{ch:v8-fluctuation-dissipation-mobility-origin-gate}

\section{Наиболее экономный кандидат}

Предыдущий гейт отделил гессиан $H$ от мобильности $M$. Самый экономный
способ связать их --- использовать кинетическую метрику, выведенную тем же
конечным следом, и стандартную градиентную идентификацию
\begin{equation}
 M\propto K^{-1}.
 \label{eq:v8-einstein-mobility-candidate}
\end{equation}
В полевом суперсвязностном носителе Тома VII уже доказано
\begin{equation}
 \Tr(\delta\mathcal D_B\,\delta\mathcal D_B)
 =3\Tr(\delta B\,\delta B^*),
 \qquad K_B=3I_{12}.
 \label{eq:v8-cross-field-trace-kinetic-metric}
\end{equation}
Поэтому на двенадцати вещественных межсекторных компонентах кандидат
\eqref{eq:v8-einstein-mobility-candidate} изотропен:
\begin{equation}
 M_B\propto\frac13I_{12},
 \qquad \kappa_{QLYR}:\kappa_{XLdR}=1:1.
 \label{eq:v8-cross-relative-mobility-pass}
\end{equation}
Это реальный частичный результат: следовая метрика не предпочитает один из
двух межсекторных мультиплетов.

\section{Почему этот результат не является общей мобильностью}

Носитель $B$ покрывает только два выбранных типа рёбер внутри своего
биклика; ещё четыре выбранных типа лежат вне него. Он не является единым
шумовым пространством связывающего, трёх калибровочных и двух межсекторных слагаемых.
Кроме того, абсолютный теплоядровой коэффициент кинетического члена в
предыдущем гейте прямо оставлен открытым. Поэтому даже
\eqref{eq:v8-cross-relative-mobility-pass} не задаёт физическую единицу
времени.

На полном шестисемейном пространстве возникают разные естественные чтения
слова ``общий''. Число компонент семейств равно
\begin{equation}
 n_r=(1,8,3,1,6,6)
 \quad\text{для}\quad
 (L,SU(3),SU(2),U(1),QLYR,XLdR).
 \label{eq:v8-noise-family-component-counts}
\end{equation}
Равная мобильность на семейство даёт веса $(1,1,1,1,1,1)$, тогда как
равная полная шумовая мощность семейства даёт
\begin{equation}
 \left(1,\frac18,\frac13,1,\frac16,\frac16\right).
 \label{eq:v8-equal-total-noise-power-weights}
\end{equation}
Нормировка по длине супероператора и по квадрату этой длины даёт ещё два
неэквивалентных вектора. Все четыре нормированных варианта примитивны и
KMS-симметричны, но максимальное расстояние между их векторами весов равно
$0.3520$, а щели лежат между $6.49\cdot10^{-5}$ и $0.02126$.

\section{Нормировка шумового кадра}

Причина неоднозначности не только физическая, но и геометрическая.
Диссипатор квадратичен по оператору скачка:
\begin{equation}
 \mathcal D[sV]=s^2\mathcal D[V],
 \qquad
 \kappa' =\frac{\kappa}{s^2}.
 \label{eq:v8-noise-frame-rescaling}
\end{equation}
Поэтому одна и та же физическая полугруппа имеет разные численные
коэффициенты в разных нормировках шумового кадра. Четыре явных
масштабирования дали нулевой остаток восстановления генератора.

Текущая вычислительная сборка нормирует каждую межсекторную вариацию по норме
Фробениуса, но использует связывающую инцидентность без той же перенормировки.
Это корректно для проверки существования процесса, однако запрещает
считать равенство семейных коэффициентов инвариантным выводом.

\begin{theorem}[Частичная межсекторная мобильность]
Общий конечный след задаёт изотропную кинетическую метрику на существующем
двенадцатимерном межсекторном носителе и при градиентной идентификации фиксирует
отношение $\kappa_{QLYR}/\kappa_{XLdR}=1$. Он не задаёт мобильности связывающего
и калибровочных семейств, абсолютного времени и канонической общей нормировки всех
операторов скачка.
\end{theorem}

\begin{nogocriterion}[Запрет равных весов без канонического шумового кадра]
Нельзя объявлять $(1,1,1,1,1,1)$ следствием общего следа, пока все скачки не
вложены в один физически типизированный модуль с одной нормой. Равенство
коэффициентов меняется при $V_r\mapsto s_rV_r$, хотя сам генератор остаётся
неизменным после $\kappa_r\mapsto\kappa_r/s_r^2$.
\end{nogocriterion}

\begin{protocol}
\begin{enumerate}
 \item следовая межсекторная метрика: \textbf{$3I_{12}$};
 \item относительная межсекторная мобильность: \textbf{$1:1$, условно получена};
 \item абсолютная межсекторная скорость: \textbf{не получена};
 \item покрытие всех шести QMS-семейств: \textbf{нет};
 \item естественных нормировок полного семейства: \textbf{4 различных};
 \item все представители примитивны и KMS-симметричны: \textbf{да};
 \item максимальный разброс нормированных весов: \textbf{$0.3520$};
 \item инвариантность генератора при смене кадра: \textbf{точно};
 \item единая шестисемейная мобильность: \textbf{не выведена};
 \item следующий гейт: \textbf{канонический шумовой кадр одного общего следа}.
\end{enumerate}
\end{protocol}

Машинный сертификат сохранён в
\path{s2t_v8_fluctuation_dissipation_mobility_origin_gate_results.json}.

\begin{thebibliography}{9}
\bibitem{v8-mobility-carlen-maas}
E.~A.~Carlen, J.~Maas,
\emph{Gradient flow and entropy inequalities for quantum Markov semigroups
with detailed balance}, J. Funct. Anal. 273 (2017) 1810--1869;
arXiv:1609.01254.
\bibitem{v8-mobility-wirth}
M.~Wirth,
\emph{The differential structure of generators of GNS-symmetric quantum
Markov semigroups}, arXiv:2207.09247.
\bibitem{v8-mobility-davies}
E.~B.~Davies,
\emph{Markovian Master Equations}, Commun. Math. Phys. 39 (1974) 91--110.
\end{thebibliography}